% GENERATED FILE. Edit canonical lessons, then run npm run generate:books.
% canonical-source-hash: fnv1a64:d420668c12c1de77
% canonical-lessons: JA-C260-futtara, JA-C260-moshi, JA-C260-moshi-tara, JA-R260-if-it-rains

\chapter{If It Rains}
\label{ch:ja-if-it-rains}
\hlchaptermodality{260}

\begin{chapteropening}
\textbf{By the end of this chapter:} \emph{I can say what will happen if something happens (ame ga futtara, ikimasen), what I will do once something happens (eki ni tsuitara), and ask what someone will do if (moshi … -tara, dō shimasu ka).}

Complete JA-R260-if-it-rains: say what you will do if it rains, and once you arrive.
\end{chapteropening}

\section[\texorpdfstring{\ja{ふったら} (futtara)}{ふったら (futtara)}]{\ja{ふったら} (futtara) — if it rains}

\label{lesson:JA-C260-futtara}



\begin{quote}
Say, plainly, \textquotedblleft{}it rained\textquotedblright{}. (\emph{ame ga futta}.) Then \textquotedblleft{}I won't go\textquotedblright{}, politely. (\emph{ikimasen}.)
\end{quote}

\subsection*{You'll want to know: -\ja{たら}, if}
\begin{quote}
\textbf{\ja{あめが} \ja{ふったら、いきません。}}
\end{quote}

\begin{quote}
\emph{ame ga futtara, ikimasen.} — \textquotedblleft{}If it rains, I won't go.\textquotedblright{}
\end{quote}

Take the plain past and add \textbf{\ja{ら}}: \textbf{\ja{ふった}} becomes \textbf{\ja{ふったら}} (\emph{futtara}), \textquotedblleft{}if it rains\textquotedblright{}. The condition comes first; what follows it comes second, and that last verb says when.

\begin{quote}
\textbf{\ja{じかんが} \ja{あったら、えいがを} \ja{みます。}}
\end{quote}

\begin{quote}
\emph{jikan ga attara, eiga o mimasu.} — \textquotedblleft{}If I have time, I will see a film.\textquotedblright{}
\end{quote}

An i-adjective does the same from its \textbf{-\ja{かった}}: \textbf{\ja{さむかった}} gives \textbf{\ja{さむかったら}} (\emph{samukattara}).

\begin{quote}
\textbf{\ja{さむかったら、うちに} \ja{います。}}
\end{quote}

\begin{quote}
\emph{samukattara, uchi ni imasu.} — \textquotedblleft{}If it is cold, I will stay home.\textquotedblright{}
\end{quote}

\subsection*{You'll want to know: -\ja{たら}, once}
When the first thing is sure to happen, \textbf{-\ja{たら}} is \textquotedblleft{}once\textquotedblright{}, \textquotedblleft{}when\textquotedblright{}:

\begin{quote}
\textbf{\ja{えきに} \ja{ついたら、きっぷを} \ja{かいます。}}
\end{quote}

\begin{quote}
\emph{eki ni tsuitara, kippu o kaimasu.} — \textquotedblleft{}Once I get to the station, I will buy a ticket.\textquotedblright{}
\end{quote}

\begin{quote}
\textbf{\ja{うちに} \ja{かえったら、ねます。}}
\end{quote}

\begin{quote}
\emph{uchi ni kaettara, nemasu.} — \textquotedblleft{}When I get home, I will go to bed.\textquotedblright{}
\end{quote}

\subsection*{Your turn}
Say these aloud:

\begin{itemize}
\raggedright
  \item \textquotedblleft{}If it rains, I won't go\textquotedblright{} — \emph{ame ga futtara, ikimasen}
  \item \textquotedblleft{}If it snows, I won't go\textquotedblright{} — \emph{yuki ga futtara, ikimasen}
  \item \textquotedblleft{}If it is cold, I won't go\textquotedblright{} — \emph{samukattara, ikimasen}
\end{itemize}

\begin{tcolorbox}[breakable,colback=teal!4,colframe=teal!35!black,title={Before you move on}]
\textquotedblleft{}If it rains\textquotedblright{}? (\textbf{\ja{ふったら}}.) \textquotedblleft{}If I have time\textquotedblright{}? (\textbf{\ja{じかんが} \ja{あったら}}.) \textquotedblleft{}Once I get home\textquotedblright{}? (\textbf{\ja{かえったら}}.) How is \textbf{-\ja{たら}} made? (The plain past, then \textbf{\ja{ら}}.)
\end{tcolorbox}
\section[\texorpdfstring{\ja{もし} (moshi)}{もし (moshi)}]{\ja{もし} (moshi) — if; supposing}

\label{lesson:JA-C260-moshi}



\begin{quote}
Before the new word: say \textquotedblleft{}if it rains, I won't go\textquotedblright{}. (\emph{ame ga futtara, ikimasen}.) Then say \textquotedblleft{}hello\textquotedblright{}, on the telephone. (\emph{moshimoshi}.)
\end{quote}

\subsection*{You'll want to know: \ja{もし}}
\textbf{\ja{もし}} — \emph{moshi} — \textquotedblleft{}if\textquotedblright{}, \textquotedblleft{}supposing\textquotedblright{}. It opens the sentence and warns the listener that a condition is coming; the \textbf{-\ja{たら}} at the end of the clause still carries the \textquotedblleft{}if\textquotedblright{}.

\begin{quote}
\textbf{\ja{もし} \ja{あめが} \ja{ふったら、いきません。}}
\end{quote}

\begin{quote}
\emph{moshi ame ga futtara, ikimasen.} — \textquotedblleft{}If it rains, I won't go.\textquotedblright{}
\end{quote}

It is not the \textbf{\ja{もし}} of \textbf{\ja{もしもし}}, the telephone \textquotedblleft{}hello\textquotedblright{}.

\begin{cousinweb}
\textbf{\ja{もし}} (\emph{moshi}) is a native Japanese word (wago), inherited rather than borrowed from Chinese. JMdict marks it as usually written in kana and glosses it \textquotedblleft{}if; in case; supposing\textquotedblright{} (JMdict).
\end{cousinweb}

\subsection*{Your turn}
Say these aloud:

\begin{itemize}
\raggedright
  \item \emph{moshi}
  \item \textquotedblleft{}If it rains, I won't go\textquotedblright{} — \emph{moshi ame ga futtara, ikimasen}
  \item \textquotedblleft{}If it snows, I won't go\textquotedblright{} — \emph{moshi yuki ga futtara, ikimasen}
\end{itemize}

\begin{tcolorbox}[breakable,colback=teal!4,colframe=teal!35!black,title={Before you move on}]
\textquotedblleft{}If\textquotedblright{}, at the start? (\textbf{\ja{もし}}, \emph{moshi}.) Which part still carries the \textquotedblleft{}if\textquotedblright{}? (The \textbf{-\ja{たら}}.)
\end{tcolorbox}
\section[\texorpdfstring{\ja{もし} … \ja{たら} (moshi … tara)}{もし … たら (moshi … tara)}]{\ja{もし} … \ja{たら} (moshi … tara) — if … then}

\label{lesson:JA-C260-moshi-tara}



\begin{quote}
Say \textquotedblleft{}if it rains, I won't go\textquotedblright{}, with \textbf{\ja{もし}}. (\emph{moshi ame ga futtara, ikimasen}.) Then \textquotedblleft{}what do you think?\textquotedblright{} (\emph{dō omoimasu ka}.)
\end{quote}

\subsection*{You'll want to know: \ja{もし} … -\ja{たら}}
\begin{quote}
\textbf{\ja{もし} \ja{はれたら、やまに} \ja{いきます。}}
\end{quote}

\begin{quote}
\emph{moshi haretara, yama ni ikimasu.} — \textquotedblleft{}If it clears up, I will go to the mountains.\textquotedblright{}
\end{quote}

The frame holds the whole condition: \textbf{\ja{もし}} opens it, \textbf{-\ja{たら}} closes it, and \textbf{\ja{、}} comes after it. Everything between is the condition:

\begin{quote}
\textbf{\ja{もし} \ja{あした} \ja{はれたら、やまに} \ja{いきます。}}
\end{quote}

\begin{quote}
\emph{moshi ashita haretara, yama ni ikimasu.} — \textquotedblleft{}If it is fine tomorrow, I will go to the mountains.\textquotedblright{}
\end{quote}

\subsection*{You'll want to know: \ja{もし} … -\ja{たら、どう} \ja{しますか}}
To ask what someone will do if, \textbf{\ja{どう} \ja{しますか}} (\emph{dō shimasu ka}), \textquotedblleft{}what will you do\textquotedblright{}, follows the frame:

\begin{quote}
\textbf{\ja{もし} \ja{あめが} \ja{ふったら、どう} \ja{しますか。}}
\end{quote}

\begin{quote}
\emph{moshi ame ga futtara, dō shimasu ka.} — \textquotedblleft{}If it rains, what will you do?\textquotedblright{}
\end{quote}

\begin{quote}
\textbf{\ja{うちで} \ja{えいがを} \ja{みます。}}
\end{quote}

\begin{quote}
\emph{uchi de eiga o mimasu.} — \textquotedblleft{}I will see a film at home.\textquotedblright{}
\end{quote}

\subsection*{Your turn}
Say these aloud:

\begin{itemize}
\raggedright
  \item \textquotedblleft{}If it is fine tomorrow, I will go to the mountains\textquotedblright{} — \emph{moshi ashita haretara, yama ni ikimasu}
  \item \textquotedblleft{}If it is fine tomorrow, I will go to the sea\textquotedblright{} — \emph{moshi ashita haretara, umi ni ikimasu}
  \item \textquotedblleft{}If it is fine tomorrow, I will go to the river\textquotedblright{} — \emph{moshi ashita haretara, kawa ni ikimasu}
\end{itemize}

\begin{tcolorbox}[breakable,colback=teal!4,colframe=teal!35!black,title={Before you move on}]
Which word opens the condition? (\textbf{\ja{もし}}.) Which closes it? (\textbf{-\ja{たら}}.) \textquotedblleft{}If it rains, what will you do?\textquotedblright{} (\textbf{\ja{もし} \ja{あめが} \ja{ふったら、どう} \ja{しますか}}.)
\end{tcolorbox}
\section[(review)]{Review — if it rains}

\label{lesson:JA-R260-if-it-rains}



\begin{quote}
Say \textquotedblleft{}if it rains, I won't go\textquotedblright{}, with \textbf{\ja{もし}}. (\emph{moshi ame ga futtara, ikimasen}.)
\end{quote}

\subsection*{Your turn}
Say these aloud:

\begin{itemize}
\raggedright
  \item If I have time, I will see a film.
  \item If it is cold, I will stay home.
  \item Once I get to the station, I will buy a ticket.
  \item If it snows, what will you do?
  \item If it is fine, I will go to the mountains.
  \item It takes about two hours to the mountains.
  \item If I lose my wallet, I will go to the police box.
\end{itemize}

\begin{tcolorbox}[breakable,colback=teal!4,colframe=teal!35!black,title={Before you move on}]
\textquotedblleft{}If it rains\textquotedblright{}? (\textbf{\ja{ふったら}}.) \textquotedblleft{}If it is cold\textquotedblright{}? (\textbf{\ja{さむかったら}}.) \textquotedblleft{}Once I get home\textquotedblright{}? (\textbf{\ja{かえったら}}.) The word that opens a condition? (\textbf{\ja{もし}}.)
\end{tcolorbox}
